\documentclass[10pt,a4paper]{amsart}
\usepackage[margin=19mm]{geometry}
\usepackage{amsmath,amssymb,mathtools}
\usepackage[scr=rsfs,cal=euler]{mathalfa}
\usepackage{xfrac}
\usepackage[unicode,hidelinks,bookmarks=false]{hyperref}
\newcommand{\ie}{i.e.\ }
\newcommand{\cf}{cf.\ }
\newcommand{\resp}{respectively\ }
\newcommand{\Subsection}{\subsection}
\providecommand{\nobreakdash}{\nobreak\hbox{-}}
\setlength{\emergencystretch}{2em}
\multlinegap0pt
\usepackage{bm}
\usepackage{bbm}
\usepackage{dsfont}
\usepackage{xspace}
\usepackage{cancel}
\usepackage{mathtools}
\usepackage{amsthm}
\makeatletter
\@ifundefined{theorem}{\newtheorem{theorem}{Theorem}}{}
\@ifundefined{lemma}{\newtheorem{lemma}[theorem]{Lemma}}{}
\makeatother
\newcommand{\col}{\colon \thinspace}
\title[A smooth family of $G\_2$-instantons over a generalised Kumme]{A smooth family of $G\_2$-instantons over a generalised Kummer construction}
\author{Dominik Gutwein}
\date{Source: arXiv:2507.00655v1 (CC BY 4.0)}
\begin{document}
\maketitle

\noindent\emph{Source and licence.} A smooth family of $G_2$-instantons over a generalised Kummer construction. Version arXiv:2507.00655v1 (CC BY 4.0), distributed under the Creative Commons Attribution 4.0 International licence, as recorded in the arXiv licence metadata for this exact version and shown on the abstract page https://arxiv.org/abs/2507.00655v1. Full rights notice: licenses/arxiv-2507.00655-CC-BY-4.0.txt.\par\smallskip

\noindent\textit{Editorial scope.} Counted unit: the Lemma whose source label is \texttt{lem: fixed-points of families} in the article (source0.tex lines 521--532, source label \texttt{lem: fixed-points of families}), delivered together with the article's own complete original proof of it (source0.tex lines 534--552, one \texttt{proof} environment of 19 lines). No mathematics is written, repaired or re-derived: statement and proof are the article's own text line for line. Required context retained uncounted: none. Every definition, display and label of those lines is retained unchanged. Omitted from this package and not redistributed: 1--520, 533--533, 553--1659. Nothing inside a retained excerpt is omitted or altered: every retained body is delivered exactly as the article's lines are written, so no omission receipt of any kind accompanies this package. Citations: the retained excerpts carry no citation command at all, so no bibliography entry of the article is transcribed and no reference is redistributed. Reference adaptation: none was needed and none was made; every reference inside the delivered unit resolves inside it, so no unresolved reference or citation is produced. Printed slips kept literal: no printed word, symbol, hypothesis or number of the counted statement or of its proof was altered. Layout and class: the article's own class is \texttt{scrartcl} with packages including inputenc, fontenc, babel, verbatim, abstract, url, hyperref, graphicx, latexsym, amsmath, amssymb, amsthm, mathtools, tikz; the wrapper is the lane's pinned \texttt{amsart} editorial wrapper at 10pt on A4, which restates the theorem-like environments the retained text needs and adds the article's own macro definitions that the retained text needs (\texttt{\textbackslash col}), with extra wrapper packages (bm, bbm, dsfont, xspace, cancel, mathtools). The delivered numbering is the unit's own. Assets: no retained span contains a figure environment, a graphics-inclusion command, a PDF-page inclusion, a table or a TikZ picture; no image file is copied into this package, the package declares no asset dependency and no image file is redistributed with it. Licence: version arXiv:2507.00655v1 of this article is distributed under the Creative Commons Attribution 4.0 International licence, as recorded in the arXiv licence metadata for this exact version and shown on the abstract page https://arxiv.org/abs/2507.00655v1. A full rights notice is delivered at licenses/arxiv-2507.00655-CC-BY-4.0.txt. Characters: every retained excerpt is pure ASCII, so the delivered engine, XeLaTeX with its default Latin Modern text font, reports no missing character.
\begin{lemma}\label{lem: fixed-points of families}
Let $(B,\Vert \cdot \Vert)$ be a Banach space and $U \coloneqq \overline{B_R(0)}\subset B$ be a closed ball of radius $R\in (0,\infty]$. Assume further that $\mathfrak{F}$ is a manifold (possibly non-compact and/or with boundary), and $E \col \mathfrak{F} \times U \to U$ is a family of contractions (i.e. $E_{\mathfrak{f}} \coloneqq E(\mathfrak{f},\cdot)$ is a contraction for every $\mathfrak{f} \in \mathfrak{F}$) that satisfies $\Vert E_{\mathfrak{f}}(b_1) - E_{\mathfrak{f}}(b_2) \Vert \leq C \Vert b_1 - b_2 \Vert$ for every $\mathfrak{f}\in \mathfrak{F}$ and $b_1,b_2\in U$ for an $\mathfrak{f}$-independent constant $C<1$. For any $\mathfrak{f}\in \mathfrak{F}$ we denote by $a_{\mathfrak{f}}\in U$ the unique fixed-point of $E_{\mathfrak{f}}$ and by $\textup{Fix}\col \mathfrak{F} \to U$ the map $\mathfrak{f} \mapsto a_{\mathfrak{f}}$. The following hold:
\begin{enumerate}
\item If $E$ is continuous (as a map $\mathfrak{F} \times U \to U$), then so is $\textup{Fix}$.
\item If $E$ is $C^k$ (as a map $\mathfrak{F} \times U \to U$), then so is $\textup{Fix}$. Furthermore, its derivative $D\textup{Fix} \col T \mathfrak{F} \to U\times B$ satisfies for every $\mathfrak{f}\in \mathfrak{F}$ and $v\in T_\mathfrak{f}\mathfrak{F}$ 
\begin{align*}
\Vert (D \textup{Fix})_\mathfrak{f}(v) \Vert &\leq (1- \Vert (\partial_BE)_{(\mathfrak{f},a_\mathfrak{f})} \Vert_{\textup{op}})^{-1} \Vert (\partial_{\mathfrak{F}} E)_{(\mathfrak{f},a_\mathfrak{f})}(v) \Vert \\
& \leq  (1- c_\mathfrak{f})^{-1} \Vert (\partial_{\mathfrak{F}} E)_{(\mathfrak{f},a_\mathfrak{f})}(v) \Vert.
\end{align*} 
Here $\partial_{\mathfrak{F}}E$ and $\partial_BE$ denote the respective partial derivatives of $E$ in $\mathfrak{F}$ and $B$ direction and $c_\mathfrak{f}<C$ denotes the optimal constant for which $ \Vert E_\mathfrak{f}(b_1) -E_\mathfrak{f}(b_2) \Vert \leq c_\mathfrak{f} \Vert b_1 - b_2 \Vert$ for every $b_1,b_2 \in U$.
\end{enumerate}
\end{lemma}

\begin{proof}
In order to prove the first point, let $\mathfrak{f}_1,\mathfrak{f}_2 \in \mathfrak{F}$. Then 
\begin{align*}
\Vert a_{\mathfrak{f}_1} -a_{\mathfrak{f}_2} \Vert & \leq \Vert E_{\mathfrak{f}_1} (a_{\mathfrak{f}_1}) -E_{\mathfrak{f}_2}(a_{\mathfrak{f}_1}) \Vert + \Vert E_{\mathfrak{f}_2}(a_{\mathfrak{f}_1}) -E_{\mathfrak{f}_2}(a_{\mathfrak{f}_2}) \Vert \\
& \leq \Vert E_{\mathfrak{f}_1}(a_{\mathfrak{f}_1}) -E_{\mathfrak{f}_2}(a_{\mathfrak{f}_1}) \Vert +C \Vert a_{\mathfrak{f}_1} -a_{\mathfrak{f}_2} \Vert 
\end{align*}
where $C<1$ is $\mathfrak{f}_2$-independent by the assumption. Absorbing the second term on the right-hand side gives 
\begin{equation} \label{eq: continuity family contractions} 
\Vert a_{\mathfrak{f}_1} -a_{\mathfrak{f}_2} \Vert \leq C^\prime \Vert E_{\mathfrak{f}_1}(a_{\mathfrak{f}_1}) -E_{\mathfrak{f}_2}(a_{\mathfrak{f}_1}) \Vert.
\end{equation} 
When $E$ is continuous, then $\mathfrak{f}_2 \to \mathfrak{f}_1$ implies therefore $a_{\mathfrak{f}_2} \to a_{\mathfrak{f}_1}$ which proves the first assertion. 

For the second point, we work locally and assume that $\mathfrak{F}$ is an open subset of $\mathbb{R}^n$ or $\mathbb{R}^n_{x_1\geq 0}$ and prove differentiability of $\textup{Fix}$ at $\mathfrak{f}_0=0$. By the differentiability of $E$ we have
\begin{align*}
a_\mathfrak{f} - a_0 &= E_\mathfrak{f}(a_\mathfrak{f}) - E_0(a_0) = (DE)_{(0,a_0)}(\mathfrak{f},a_{\mathfrak{f}}-a_0) + o (\vert \mathfrak{f} \vert +\Vert a_{\mathfrak{f}}-a_0 \Vert) \\
&= (\partial_B E)_{(0,a_0)}(a_\mathfrak{f} - a_0) + (\partial_{\mathfrak{F}} E)_{(0,a_0)} \cdot \mathfrak{f} + o (\vert \mathfrak{f} \vert+\Vert a_\mathfrak{f} -a_0 \Vert)
\end{align*}
where $DE$ denotes the (total) differential of $E$ and $\partial_{\mathfrak{F}}E$ and $\partial_BE$ the partial derivatives in $\mathfrak{F}$ and $B$ direction. Since $\Vert E_0(b) - E_0(a_0) \Vert \leq c_0 \Vert b - a_0 \Vert$ for every $b \in U$ with $c_0<1$, we obtain $\Vert (\partial_BE)_{(0,a_0)} \Vert_{\textup{op}} = \Vert (DE_0)_{a_0} \Vert_{\textup{op}} \leq c_0 <1$. The operator $(\textup{Id} - (\partial_BE)_{(0,a_0)})$ is therefore invertible via the Neumann series of $(\partial_BE)_{(0,a_0)}$. This gives \[ a_\mathfrak{f} - a_0 = (\textup{Id} - (\partial_BE)_{(0,a_0)} )^{-1} \big((\partial_{\mathfrak{F}} E)_{(0,a_0)} \cdot \mathfrak{f} + o (\vert \mathfrak{f} \vert+\Vert a_\mathfrak{f} -a_0 \Vert)\big). \] The differentiability of $E$ together with~\eqref{eq: continuity family contractions} imply that $o(\vert \mathfrak{f}\vert + \Vert a_\mathfrak{f} -a_0 \Vert)=o(\vert \mathfrak{f} \vert)$. The previous formula for $a_\mathfrak{f}-a_0$ shows therefore that $\textup{Fix}$ is differentiable at $\mathfrak{f}=0$ with derivative \[ (D\textup{Fix})_\mathfrak{f} = (\textup{Id} - (\partial_BE)_{(\mathfrak{f},a_\mathfrak{f})} )^{-1} \big( (\partial_{\mathfrak{F}} E)_{(\mathfrak{f},a_\mathfrak{f})}(\ \cdot\ )\big). \] Since inversion and concatenation are smooth operations, $E \in C^k$ implies via bootstrapping that $\textup{Fix}\in C^k$.
\end{proof}

\end{document}
